\section{Measurement methodology}
\label{sec:method}

\subsection{Stack and instrumentation}
\label{sec:stack}

All measurements run on TP4 GB200 fleets of 4, 8, or 16 replicas
serving Qwen3-Coder-480B-A35B-Instruct-FP8 under vLLM with block
size 256; each replica holds a KV pool of $6486 \times 256$-token blocks
(13.28M tokens at 8 replicas) \evid{experiment-env.md}. Requests reach
the fleet through an Envoy front end and an Endpoint Picker (EPP)
router. The router plugin configuration is a first-class
experimental factor, selected per run among three configurations: a
precise prefix-cache configuration (KV-event token-based affinity;
the baseline for the main series), an approximate-prefix
configuration with no KV-event pipeline, and a load-only
configuration (queue and KV-utilization scorers, no affinity
signal) \evid{experiment-env.md}.

Server state is scraped from per-replica metrics endpoints at 5\,s
cadence, and fleet aggregates use replica endpoints only (the
service endpoint round-robins replicas) \evid{C34}. Hit rate $h$ is
computed as windowed counter differences of
\texttt{vllm:prompt\_tokens\_cached} over \texttt{vllm:prompt\_tokens}; the
\texttt{prefix\_cache\_queries}/\texttt{hits} pair is not used because it re-counts
under pressure \evid{experiment-env.md}. On tiered fleets, the tier hit
fraction $\hext$ is the external (host-tier) hit-token fraction,
\texttt{vllm:external\_prefix\_cache\_hits} over \texttt{external\_prefix\_cache\_queries}
(token counters), and the restore share of hits is the fraction of
hit tokens served from the tier \evid{overlay-findings.md}. Time series
are aggregated into 5-minute windows anchored to wall clock so
concurrent runs align \evid{analyze\_run\_windows.py}. One gauge-semantics
caveat applies throughout: the \texttt{vllm:num\_requests\_running} gauge
excludes requests in \texttt{WAITING\_FOR\_REMOTE\_KVS}, so server-side running
counts do not equal client-side load during restore phases \evid{R9}.

\subsection{Workload}
\label{sec:workload}

The primary corpus is a 393-session agentic coding replay set
(68,266 requests, mean $\sim$174 requests per session), collected from
real coding-agent traffic and replayed with per-instance prefix
salting so replay instances share no cache (unsalted replays inflate
$h$) \evid{experiment-env.md}. The corpus exhibits the agentic structure of
Section~\ref{sec:intro}: $\sim$52k-token opening system prompts, few-thousand-token
per-turn increments, growth to a 249k compaction ceiling, bimodal
think times (78\% of turns within seconds; a long human-pace tail),
and subagent fan-out in 44\% of sessions. Inter-request idle gaps are
capped at 10.5\,s, the corpus p90, making the replay a machine-pace
workload; the cap is an experimental control for the reuse-gap
distribution $F$ (Sections \ref{sec:retention}, \ref{sec:implications}). Surge traffic replays a disjoint
232-session corpus so burst and base load share no session identity.

Arrival semantics determine whether collapse is reachable, so we
state them precisely. Open-loop experiments use Poisson session
arrivals at rate $\ls$ with concurrency as an admission ceiling;
requests within a session remain completion-ordered. Closed-loop
experiments meter a fixed credit of in-flight work; the measured
credit multiplier is $\sim$1.45$\times$ configured concurrency on this corpus,
so all load-axis statements use achieved in-flight \evid{R8}.
Request-metered rate mode is excluded from collapse experiments: it
re-issues work independently of completions, a built-in retry storm,
whereas completion-coupled session arrivals defer rather than storm.
Open-loop collapse experiments therefore pin demand at the session
level \evid{R7}\evid{C33}.

\subsection{Perturbation protocol}
\label{sec:protocol}

Each collapse experiment resets every replica's prefix cache, runs
base load at $\ls$, and at $t \sim +62$\,min launches a surge overlay
without cache reset: 0.25 sessions/s for 2700\,s (the standard surge;
a 900\,s variant probes the duration axis), with cancellation at
$t \sim +107$\,min. The surge generator cancels its entire session
population on exit, so all post-cancellation dynamics are sustained
by the base load alone. Horizons run 150--300\,min; Section~\ref{sec:horizon} shows
why the long horizon is necessary. Closed-loop states are held 65
min before being read as stationary \evid{R8}.

\begin{figure*}[t]
\centering
\includegraphics[width=\textwidth]{fig_protocol.png}
\caption{Perturbation protocol timeline. Prefix-cache reset on all
replicas at $t{=}0$; base Poisson session arrivals at $\ls$ sustained
to the end of the 150--300\,min horizon; eviction-scale surge overlay
(0.25 sessions/s $\times$ 2700\,s, disjoint corpus, no cache reset) over
[62, 107)\,min, with the surge population cancelled at $t{=}107$ so the
post-cancellation drain-versus-catch-up race over [107, $\sim$120)\,min is
base-sustained. The warm state is read over [30, 60]\,min, the drain
depth $D$ (Section~\ref{sec:separatrix}) over [115, 120)\,min, and end-state outcomes
are classified over [270, 295]\,min (cold if mean $h < 0.15$).}
\label{fig:protocol}
\end{figure*}

\subsection{Comparability and run-to-run variance}
\label{sec:comparability}

Five factors are recorded per run and held fixed within any
comparison: router configuration, corpus, gap cap, salt, and
duration \evid{C34}; run identifiers and per-run configurations are
recorded in the artifact ledger. An early reference series that ran
under different router configurations and unsalted generators is
excluded from all fits. Realized request rate varies $\sim\pm0.4$
requests/s between runs at fixed $\ls$ and does not order
outcomes in the series where this was checked \evid{R2}\evid{C04};
consequently every probabilistic statement in this paper is an
outcome count over repeated runs, never a response inferred from a
single run. Runs are treated as statistically independent: every run
starts from a cache reset, and a namespace-split check across the
two shard environments supports run-to-run variance over a fixed
environment effect. Two dependence structures remain and are carried
as confounds (Section~\ref{sec:limitations}): shard runs executed pairwise in shared
windows, and two same-configuration replicate pairs across days.

Two artifact-existence rules make absent data auditable rather than
silently missing: load-generator artifacts exist only after
end-of-run export (one hung export in $\sim$40 runs through the
mid-campaign window, zero in the final 19-run window), and capacity
reclaim destroys un-scaled fleets, so an unconditional scale-down is
scheduled at window start \evid{R12}\evid{C35}.

\subsection{Evidence grading}
\label{sec:grading}

Every claim in this paper carries one of four grades, and the full
register appears in Appendix~\ref{sec:claims}. \emph{Demonstrated}: direct measurement,
replicated or internally cross-checked; stated as fact with $n$
reported. \emph{Outcome frequency ($n$)}: a tally over $n$ repeated runs;
stated as the tally, never as a fitted probability. \emph{Model-supported}:
holds in the calibrated fluid/DES models; attributed to the model
with its validation scope. \emph{Open}: a documented failure or unresolved
question, reported as such. Single-run observations are flagged
inline ($n=1$) and never support a probability, threshold, or trend on
their own. We regard the grading itself as methodology: metastable
phenomena produce run-to-run outcome variance at fixed operating
points (Section~\ref{sec:band}), and ungraded reporting of single runs is how
such phenomena get mis-measured.
